\section{The far region and simultaneous closure}\label{sec:far}

The compact estimates leave two quantities to control: radiation at large
radius and the acceleration of the modulated Kerr parameters.  We now prove
the far estimate with a radius-small coefficient on both quantities, and
then close it simultaneously with the compact and sliced estimates.
The far equation has no conic pseudodifferential modification.  We first
verify its target symbols, differentiated weak-null structure, and
initial and overlap bounds.  The energy and product estimates in
Appendices~\ref{fe:section} and~\ref{fs:section} then apply to this
equation.  Their role is to turn the structural bounds into a
radius-small far estimate; the final subsection here closes its
coupling to the compact region.

Throughout we work on an arbitrary smooth finite bootstrap slab.  All
constants are independent of its terminal time.  The compact norm
$\mathsf K_R$, the acceleration norm
$\mathsf S_a=\|a\|_{L^2_A}$, the total bootstrap size $\mathsf N$, and the
base input size $e_{\mathrm{in}}$ have the meanings fixed in the compact
argument.  Here $a$ in $\mathsf S_a$ denotes the modulation acceleration,
not the fixed Kerr rotation parameter.  A constant $C_{\mathrm{fix}}$ may
depend on every scale chosen before the final data smallness, but not on
the lifespan.

\subsection{Emitter coordinates and the three far norms}

Let $\eta$ denote the Minkowski metric in the lab coordinates $(t,z)$.
Write
\[
 \begin{gathered}
 C_e(q)=-\lambda(q),\qquad s_e=C_e',\qquad s_s=-v_s(t,x),\\
 t=q+r,\qquad z=C_e(q)+rn,\qquad |n|=1,
 \end{gathered}
\]
and set $k=1-s_e\cdot n$.  In this section $r$ is emitter radius; the
small speed bound makes $k$ uniformly close to one and makes $r$, $|x|$,
and the co-moving radius $|y|$ comparable.  Put
\[
 \begin{gathered}
 T_l=\partial_{t,\mathrm{lab}},\qquad
 L=T_l+n\cdot\partial_z=\partial_{r|q,n},\\
 S=rL,\qquad Z=(S,\Omega),\qquad D=T_l+s_s\cdot\partial_z,
 \end{gathered}
\]
where $\Omega$ denotes the emitter sphere rotations.  All tensor
components in this section are lab components, and we abbreviate
\[
 h=h_{\mathrm{lab}},\qquad \psi=Dh,\qquad H^{ij}=g^{ij}-\eta^{ij}.
\]
The covector $\ell_n=(1,-n)$ is used only for null contractions.  Thus
$H_{\ell_n\ell_n}=H^{ij}\ell_{n,i}\ell_{n,j}$.  Let $\Pi(n)$ be the
projection of a covariant symmetric tensor onto its tracefree spatial
screen part orthogonal to $n$.

The elementary derivative conversion is
\begin{equation}\label{far:derivative-conversion}
 \partial_i=\ell_{n,i}T_l+r^{-1}b_iZ
 =\ell_{n,i}k^{-1}\partial_q
  +r^{-1}(b_i-\ell_{n,i}k^{-1}s_e^db_d)Z.
\end{equation}
Here $b_t=0$, while $b_dZ$ is $n_dS$ plus the corresponding tangential
rotation combination.  We always commute lab time derivatives before
the $Z$ derivatives.  Choose the starting radius $R_0$ beyond all fixed
transitions and all compact kernels.  The fixed-factor starting overlap
is then physically nontrapping.  We require, in addition, that radii as
small as $cR_0^{0.45}$ lie in the lab differential and symbol regime;
these smaller radii occur only in the spacelike-anchor argument below.

We use the three levels
\begin{equation}\label{far:levels}
\begin{array}{c|ccc}
 i&0&H&P\\ \hline
 U_i&h&h&\psi\\
 N_i&9&8&7\\
 A_i&0&A_H=1-\eta_*&A=1.36\\
 p_i&0.64&1-\eta_*/6&1.62
\end{array}
\end{equation}
with a fixed $0<\eta_*\le 10^{-8}$.  For the far weights define
\[
 \delta_f=\eta_*/100,\qquad \nu=10\delta_f,
 \qquad e=0.1\delta_f,\qquad d_*=200\eta_*.
\]
These quantities are distinct from the slack in the compact
differential-current construction.  With $\langle s\rangle=(1+s^2)^{1/2}$,
put
\begin{equation}\label{far:weights}
 \begin{gathered}
 \tau_f=1+t/r,\qquad \mu=\langle q/r^{\delta_f}\rangle,
 \qquad l_f=\langle q/r^\nu\rangle,\\
 m_f=(r^{-1}+l_f^{-1})^{-1},\qquad
 W_i=r^{p_i-A_i}(t+10r)^{A_i},\\
 w_{H,j}=w_H=r^{-e}\mu^{A_H},\qquad
 w_{P,j}=w_Hl_f^{A-A_H}m_f^{-(7-j)d_*},\\
 w_{0,j}=r^{-0.02}m_-^{jd_*}m_+^{-(9-j)d_*}.
 \end{gathered}
\end{equation}
Here $m_-$ is a smooth monotone function of $-q/r^\nu$, comparable to
$1+(-q/r^\nu)_+$, constant when its argument is nonpositive, and with
bounded logarithmic-symbol derivatives.  The weight $m_+$ is the
positive-$q$ analogue capped at $r$ by harmonic addition.

For each word $I$ in $Z$ and each $j$ with $|I|+j\le N_i$, define
$\chi=Z^I(rT_l^jU_i)$.  Write $G_1\chi=r^{-1}Z\chi$ and let
$B_1\chi$ denote its full first derivative.  The bulk measure will always
be $dt\,dr\,d\omega/r$.  The preliminary far norm is the sum, over these
indices and tensor components, of
\begin{equation}\label{far:basic-norm}
 \begin{split}
 &\|W_i^{1/2}(G_1\chi,\chi/r)\|_{\mathrm{bulk}}
   +\|w_{i,j}^{1/2}B_1\chi\|_{\mathrm{bulk}}\\
 &\hspace{8mm}
   +\sup_t\|(W_i^{1/2}G_1\chi,w_{i,j}^{1/2}B_1\chi)
                                      \|_{L^2(dr\,d\omega)}.
 \end{split}
\end{equation}
We complete this norm by corrected-current terms in
\eqref{far:corrected-norm} below and denote the result by $\mathsf F$.
The weights are those of \Cref{fe:F1}.  On $t\gg r$ they satisfy
\begin{equation}\label{far:timelike-weights}
 w_{i,j}\gtrsim\langle t\rangle^{A_i}r^{-0.03},
 \qquad W_i\gtrsim t^{A_i}r^{p_i-A_i}.
\end{equation}
Differentiating $\chi$ supplies the bad-derivative controls used in the
sliced argument; the extra lower-order terms are controlled by the good
entries.  Thus \eqref{far:basic-norm} contains those earlier far inputs,
with any necessary smaller positive radial slack.

\subsection{Verification of the lab equation}

The target metric is denoted by $B$, its wave compensator by $\Theta$,
and its pure equation defect by $f_B$.  Enlarge the positive acceleration
majorant at comparable radii when necessary, and write
\[
 \widetilde a(t,r)=Cr^{-1}\int_{t-Cr}^t|a(\tau)|\,d\tau.
\]
The histories are zero before startup.  The particular constant $C$ in
this notation can be increased finitely many times.

\begin{lemma}[Far symbol and weak-null structure]\label{far:structure}
For the arbitrary fixed subextremal Kerr family constructed above, the
lab target satisfies
\begin{equation}\label{far:target-symbols}
 \begin{gathered}
 B-\eta=O(r^{-1}),\qquad DB=O(r^{-1}\widetilde a),\qquad
 \partial^ds_s=O(r^{1-d}\widetilde a)\quad(d\ge1),\\
 \Theta=O(r^{-1}\widetilde a),\qquad
 f_B=O(r^{-2}\widetilde a),\qquad Df_B=O(r^{-3}\widetilde a),\\
 \qquad r\widetilde a\lesssim\mathsf N\tau_f^{-A/2}.
 \end{gathered}
\end{equation}
These are symbol estimates at every fixed order needed here, with one
inverse power per additional lab derivative.  Their constants may depend
on the fixed family but not on either matching radius.  The equation is
\begin{equation}\label{far:wave-structure}
 \begin{gathered}
 g^{ij}\partial_{ij}h=f_h,\\
 \begin{aligned}
 f_h={}&f_B+O(r^{-2}\partial h+r^{-3}h)
             +Q_\eta(\partial h,\partial h)\\
       &+O((r^{-1}+|h|)|\partial h|^2),
 \end{aligned}
 \end{gathered}
\end{equation}
up to fixed harmless changes in $f_B$.  The first- and zeroth-order
background coefficients have a relative $O(\widetilde a)$ gain under
$D$.  The flat linear wave-gauge contraction of $h$ is
\begin{equation}\label{far:flat-gauge}
 O((r^{-1}+|h|)\partial h)+O(r^{-2}h)
                              +O(r^{-1}\widetilde a).
\end{equation}
The estimates allow smooth composition in the fields and distributed
product differentiation.  The quadratic form $Q_\eta$ has the constrained
weak-null polarization described below.
\end{lemma}

\begin{proof}
The far background in the formal lab coframe has the form
$\eta+\mathfrak h(p_s(t,x),x)$, where, at independent family argument,
\[
 \partial_p^b\partial_x^\alpha\mathfrak h
                      =O(|x|^{-1-|\alpha|}).
\]
This statement includes the coframe change and holds for every member
of the chosen fixed Kerr parameter family.  Implicit differentiation of
the smoothed coordinates gives
\[
 \partial_zx=S_x=(1-\partial_x\lambda_s)^{-1},\qquad
 T_lx=S_xv_s,\qquad S_x-I=O(r\widetilde a).
\]
The first derivatives of $S_x$ and $v_s$ are $O(\widetilde a)$, with the
corresponding differentiated symbol bounds.  In particular $Dx=0$.
Positive-delay smoothing and its signed moments imply
$\partial_x\lambda_s=O(r\widetilde a)$.  The affine flat leading term
cancels exactly in lab components.  Comparing at a point with an exact
constant-parameter Kerr metric whose affine worldline agrees there gives
the first- and second-jet errors
$O(r^{-1}\widetilde a)$ and $O(r^{-2}\widetilde a)$.
Commuting $D$ past a lab derivative introduces a derivative of $s_s$;
this proves the differentiated bounds for $B$ and its defect.

For clarity, the compensator uses the stationary profiles with orders
$H_{\mathrm{mode}},H^1_{\mathrm{mode}},V_i-e_i=O(r^{-1})$ and
$J=O(r^{-2})$ from the stationary construction.  We append the subscript
here only to distinguish the mode profiles from $H^{ij}$.
Its local gauge image consists of first-order
commutators against the smoothed mass and velocity parameters, terms
containing $J$ and $\partial_x\lambda_s$, and the gauge derivative of
$(V_i-e_i)^\flat\odot d_x\lambda_s^i$.  When only one derivative falls
on $\lambda_s$ in the last term, a derivative of the profile, or a
zeroth-order connection coefficient, supplies the extra inverse radius.
The chain identities cancel the free position and velocity terms.
Thus neither an unsmoothed acceleration nor a growing free position
occurs in a lab symbol coefficient.  Filtered unsmoothed accelerations
and the finitely many low jets needed here are bounded by
$C_{\mathrm{fix}}\mathsf N\langle q\rangle^{-A/2}$ and vanish before
startup.  These observations prove \eqref{far:target-symbols}.

Constraint propagation gives the exact lab identity
\[
 \Gamma_\mu(g)=g_{\mu\sigma}g^{ab}\Gamma(B)^\sigma_{ab}+\Theta_\mu,
\]
where $\Gamma_\mu(g)$ is the lowered wave contraction.  Expanding this
identity gives \eqref{far:flat-gauge}.  Expanding vacuum Ricci in this
gauge gives \eqref{far:wave-structure}.  The added nonlocal gauge terms
vanish in the far region; only the differential wave compensator remains.
The coefficients of the linear terms use first and second jets of $B$,
their products, and compensator jets, so their claimed $D$ gain follows
from \eqref{far:target-symbols} and $[D,\partial]$.

For a symmetric polarization $p$ satisfying
\begin{equation}\label{far:polarization-constraint}
 \ell_n^\alpha\bigl(p_{\alpha\beta}
                  -\tfrac12\eta_{\alpha\beta}\operatorname{tr}_\eta p
              \bigr)=0,
\end{equation}
the flat Christoffel identity of \Cref{fe:weak-null} says that
$Q_\eta(\ell_n\otimes p,\ell_n\otimes p)$ is a constant multiple of
$\ell_n\otimes\ell_n|\Pi(n)p|^2$; its polarized version holds as well.
Indeed the constraint eliminates $p_{LL}$, $p_{LA}$, and the screen
trace in a flat null frame.  The linear map to the four contractions in
\eqref{far:polarization-constraint} is surjective on $\ker\Pi$.
Consequently there is a uniformly smooth projection to the constrained
polarizations that leaves $\Pi p$ fixed and changes $p$ by bounded
multiples of those four contractions; this is
\Cref{fe:gauge-projection}.  Apply the differentiated identity
\eqref{far:flat-gauge} to $p=T_l^{j+1}h$ before taking this projection.
After multiplying a first jet by $r$, every error contains a good
derivative, a vanishing coefficient times an ordinary scaled jet, a
symbol-suppressed field, or $O(\widetilde a)$.  Distributed lab-time
differentiations retain exactly these placements.  This verifies the
differentiable weak-null structure, not just its undifferentiated
algebraic identity.
\end{proof}

\begin{lemma}[Initial and overlap inputs]\label{far:boundary-input}
The initial, current-correction, and field overlap costs at the
fixed-factor starting annulus in the far energy estimates are bounded by
\begin{equation}\label{far:overlap-bound}
 C(R_0)\mathsf K_R+C_{\mathrm{fix}}(e_{\mathrm{in}}+\mathsf N^2).
\end{equation}
The initial input requires only the base metric $H^{10}$ and second-form
$H^9$ bounds and their order $-1,-2$ end-symbol bounds.  It imposes no
condition on leading tail coefficients and no smallness on higher
seminorms.  Pure acceleration forcing is estimated separately, including
on this annulus, by the direct averaging estimate
\eqref{fs:average-bounds} in Appendix~\ref{fs:section}.
\end{lemma}

\begin{proof}
Initially the lapse and shift have background values and their
derivatives are fixed by the gauge.  At infinity startup has zero
parameters and no kernel terms.  The wave equation supplies mixed jets
through order ten, with one inverse radius per derivative.  The incoming
weighted integrals converge because $p_0,p_H<1$ for $\chi(h)$, whereas
$\psi$ has an additional inverse power.  Reduced local propagation gives
finiteness of all higher symbol orders on bounded time intervals without
making their constants small.

On a fixed-factor overlap at $R_0$, all unscaled mixed lab jets of $U_i$
through order $N_i+1$ have full compact bulk and endpoint control, with
squared time weight $\langle t\rangle^{A_i}$.  The component transformation
from $h_{\mathrm{co}}$ to $h_{\mathrm{lab}}$ depends pointwise on $v(t)$.
Under this transformation $\psi$ equals $u$ plus terms of the forms
$(v(t)-v_s)\partial h_{\mathrm{co}}$ and $v'h_{\mathrm{co}}$ with smooth
bounded component factors.  Their coefficients carry the strong time
weight by the fixed-interval acceleration averages and parameter jet
bounds.  Inverse-coordinate and differentiated-frame costs are bounded
on this fixed overlap, and $\partial^{\le9}h_{\mathrm{co}}$ has at least
the weak weighted bound.  These terms therefore contribute products,
not a radius-growing linear coefficient on $\mathsf S_a$.
The emitter commutations have bounded conversion coefficients there.
The weights $W_i,w_{i,j}$, and the factor $1+M_f$ defined below cost at
most the indicated time weights, up to $R_0$-dependent constants; for
$1+M_f$ the cost is bounded.  This proves \eqref{far:overlap-bound} and
verifies the boundary hypotheses of \Cref{fs:claim-boundary}.
\end{proof}

\subsection{Point bounds and corrected energy currents}

We record the precise point information entering the source estimates.
It also supplies the timelike point bounds used in the sliced comparison.
For the following assertions, $N\le C_{\mathrm{fix}}\mathsf N$ denotes
the small factor used in product estimates, not a new independent norm.
The finite-slab implications of \Cref{fe:spacelike-anchors},
\Cref{fe:point-estimates}, \Cref{fe:null-projection},
\Cref{fe:tight-projection}, and \Cref{fe:F11} give
\begin{align}
 \|\chi(h)\|_{L^2_\omega}
   &\lesssim Nr^{2\eta_*}\tau_f^{-A_H/2}
       &&\text{through order eight},\label{far:point-h}\\
 \|\chi(h)\|_{L^2_\omega}
   &\lesssim Nr^{0.52}&&\text{at order nine},\notag\\
 \|B_1\chi(h)\|_{L^2_\omega}+\|\chi(\psi)\|_{L^2_\omega}
   &\lesssim Nr^{e/2}\mu^{-A_H/2}
       &&\text{through order seven}.\label{far:point-first}
\end{align}
The first-derivative version for $\psi$ holds one order lower.  Two spare
rotations give the corresponding sphere supremum bounds.  Contract lab
time jets before applying $Z$ to $H_{\ell_n\ell_n}$.  Through order eight
their angular $L^2$ bounds are
\begin{equation}\label{far:projected-bounds}
 O(r^{-1}+Nr^{-1+2\eta_*}\tau_f^{-A_H/2}),
 \qquad O(r^{-2}+Nr^{-2+8\eta_*}\tau_f^{-A_H/2})
\end{equation}
in general and with at least one lab time, respectively.  The second
bound also holds after $\partial_q$ at orders through seven.  Without a
lab time, $r$ times the coefficient is $O(1)$ through order eight on
$|q|\lesssim r^{0.45}$.  At order nine it is bounded by
\begin{equation}\label{far:top-projected}
 O\bigl(\log r+b^{1/2}r^{-p_0/2}+br^{-1+8\eta_*}\bigr),
 \qquad b=1+\min(|q|,r);
\end{equation}
with a lab time the bound for $r$ times the order-nine coefficient
improves to $O(r^{-0.46})$.  Finally, the lowest radiative jets satisfy
\begin{equation}\label{far:news}
 |r\Pi T_lh|+|r\Pi T_l^2h|+|r\Pi T_l\psi|\lesssim N.
\end{equation}

There is a radius-uniform point in applying these implications.
The negative-$q$ spacelike anchor is proved on an auxiliary exterior
\begin{equation}\label{far:anchor-domain}
 \begin{gathered}
 |z|\ge t+cQ+C_1\int_0^t((u+Q)Q)^{-1/2}\,du,\\
 0\le t\lesssim R_f,\qquad Q\gtrsim R_f^{0.45}.
 \end{gathered}
\end{equation}
Here $R_f$ is a dyadic radius.  This domain has $q<0$; zero startup
histories give $\lambda(q)=0$ and hence $r=|z|$.  Its lower boundary may
lie below $R_0$, but our
choice of $R_0$ places all of it in the differential lab end.  Hence the
same equation and initial symbol bounds apply there.  Regard the
histories as fixed in this domain estimate.  The pure source after up to
nine Minkowski commutations costs
$O(\mathsf S_a|z|^{-3})$ by the same acceleration majorant.
Thus the moving-boundary energy bootstrap of
\Cref{fe:spacelike-anchors} uses no radius-dependent replacement by
$\mathsf K_R$.  Constants acting linearly on far quantities remain
bounded independently of $R_0$; nonlinear smallness can be imposed after
the radii are fixed.  The news estimate \eqref{far:news} follows from the
emitter equations on $rT_l^jh$, $j=0,1$: projection kills the leading
quadratic by Lemma~\ref{far:structure}.  No point smallness in a deep
compact region is used, except on the prescribed overlaps.

To describe the currents, set
\[
 P_g\chi=krg^{ij}\partial_{ij}(\chi/r),\qquad
 P[J]\chi=krJ^{ij}\partial_{ij}(\chi/r).
\]
The exact flat emitter operator and the principal perturbation identity are
\begin{align}
 P_\eta={}&-2\partial_q\partial_r+r^{-2}\bigl[
 (1+s_e\cdot n)(S^2-S)-2(s_e\cdot n)(S-1)\notag\\
 &\hspace{40mm}+2s_e^A\nabla_A(S-1)+k\Delta_{\mathbb S^2}\bigr],
                                                        \label{far:flat-operator}\\
 P[J]\chi={}&J_{\ell_n\ell_n}\partial_q(k^{-1}\partial_q\chi)
     +O\bigl((J/r)(1,Z)\partial_q\chi
                       +(J/r^2)(1,Z,Z^2)\chi\notag\\
 &\hspace{42mm}+(a(q)/r)J_{\ell_n\ell_n}(1,Z)\chi\bigr).
                                                        \label{far:perturbation}
\end{align}
The $O$ notation here denotes finite sums with bounded smooth
emitter-frame coefficients.  These identities are \Cref{fe:F4} and
\Cref{fe:F5}; they are unchanged by the choice of Kerr spin.


Let $X_g,\bar X_g$ be the true future outgoing and incoming null normals,
normalized by $dt=1$, let $G=-g(X_g,\bar X_g)$, and let $\alpha$ be the
$\bar X_g$ coefficient of $\partial_q$ modulo $X_g$ and the screen.
Define
\begin{align}
 d_g&=|H_{\ell_n\ell_n}|+|H|^2,\qquad
 M_f=r^{0.72}\mu^{-A_H}j_0(q/r),\notag\\
 \mathfrak u&=q+\int_{R_0/2}^r
       \left[\rho^{\delta_f-1}
         +\rho^{20\eta_*-1}j(|q|/\rho^{0.45})\right]d\rho.
                                                        \label{far:caps}
\end{align}
The smooth cutoffs satisfy $j_0=1$ on $[-1,1]$, $j_0=0$ off $[-2,2]$,
$j=0$ below one, and $j=1$ above two.  The tight projection bounds make
the caps $\mathfrak u=u_0$ spacelike when $R_0$ is large.

For each commuted equation the outgoing source is split as
\begin{equation}\label{far:source-split}
 P_g\chi=\lambda_1L\chi/r+f_{\mathrm{flat}}
                       +\partial_q\mathcal P+f_{\mathrm r}+f_{\mathrm{mix}}.
\end{equation}
Here $\lambda_1$ is bounded and nonnegative; it and the triangular
$f_{\mathrm{flat}}$ arise from flat dilation and rotation commutators.
Write $f_{\mathrm{raw}}$ for the unsplit remaining source after these
flat terms.  The correction $\mathcal P$ is formed term by term in the commuted
lab equation, keeping each algebraic sign.  In the following list,
$\chi'=Z^{I'}(rT_l^{j'}U)$ is the commuted input field,
$U\in\{h,\psi\}$; a coefficient $Z^dT_l^mH$ has count $d+m$.
For the smooth-acceleration terms, $C$ denotes the twice-null
contraction of an ordered derivative of $2\eta^{ij}\partial_i s_s^d$,
the coefficient of $\partial_{jd}h$ in the strong equation.
There are three types:
\begin{enumerate}[label=(\roman*)]
\item A double-bad principal term
$C_h\partial_q(c_0\partial_q\chi')$ contributes
$C_hc_0\partial_q\chi'$ when the coefficient count is at most five,
or at every coefficient count at level $P$.
Here $C_h=Z^d(T_l^mH)_{\ell_n\ell_n}$ and
$c_0=c_0(n,s_e(q))$ is the smooth emitter coefficient.
The smooth-acceleration double-bad terms in the $D$ equation use the
same rule with their coefficient in place of $C_h$.
\item A frame term $C(\partial_qc_2)D_Z\chi'/r$ in that equation
contributes $Cc_2D_Z\chi'/r$, where $c_2=c_2(n,s_e(q))$ is smooth
and bounded, and $D_Z$ has $Z$ order at most one and is independent of $q$.
\item Put $F_m=r\Pi T_l^{m+1}h$.  Each leading strong quadratic
pairing contributes
\begin{equation}\label{far:leading-correction}
 \mathcal P_*=Z^{b'}\!\left[(1-s_s\cdot n)r^{-1}\ell_n\otimes\ell_n
                       \mathcal C_{mm'}(F_m,F_{m'})\right],
\end{equation}
where $\mathcal C_{mm'}$ is its actual polarized screen pairing.
\end{enumerate}
The ordered expansions and these choices are given in
\eqref{fs:G5}, \eqref{fs:G7}, and \eqref{fs:G10}--\eqref{fs:G12}.
The raw source keeps all three types unsplit for the lab-time multiplier.
With their signed sum $\mathcal P$, define
\[
 \xi_{\mathrm{curr}}
   =d\chi+\frac{G\alpha}{2k}\mathcal P
                   \frac{-g(\bar X_g,\cdot)}{G}.
\]
We now complete $\mathsf F$ by adding, for each equation,
\begin{equation}\label{far:corrected-norm}
 \|(W_i(1+M_f))^{1/2}\xi_{\mathrm{curr},X_g}\|_{\mathrm{bulk}}
 +\sup_{u_0}\|W_i^{1/2}\xi_{\mathrm{curr},X_g}
                   \|_{L^2(dr\,d\omega;\,\mathfrak u=u_0)}.
\end{equation}
The source costs of \Cref{fe:energy-reduction} are
\begin{equation}\label{far:source-costs}
 \begin{split}
 &\|W_i^{1/2}(\mathcal P,Z\mathcal P,rd_g\partial_q\mathcal P)
                                                    \|_{\mathrm{bulk}}
       +\sup_t\|W_i^{1/2}\mathcal P\|_{L^2_{r,\omega}}\\
 &\quad+\|(W_i/(1+M_f))^{1/2}rf_{\mathrm r}\|_{\mathrm{bulk}}
       +\|w_{i,j}^{1/2}rf_{\mathrm{raw}}\|_{\mathrm{bulk}}.
 \end{split}
\end{equation}
There is also a mixed cap charge: on $t\lesssim r$, $r\sim R_f$, a bound
$\|rf_{\mathrm{mix}}\|_{L^2_\omega}
\le l_{R_f}(\mathfrak u)\mathcal T(t,r)$ is charged by
\begin{equation}\label{far:mixed-charge}
 \sum_{R_f}R_f^{(p_i-1)/2}\|l_{R_f}\|_{L^2(du_0)}
                            \sup_t\|\mathcal T(t,\cdot)\|_{L^2_r}.
\end{equation}

The geometric assumptions of the energy reduction follow from the
point bounds already recorded:
\[
 X_gq=O(d_g),\qquad
 d_g\lesssim r^{-1}\ \text{if }|q|\le2r^{0.45},\qquad
 d_g\lesssim r^{-1+2\eta_*}\ \text{otherwise}.
\]
Good derivatives of lab metric components, and an unrestricted first
derivative when both metric slots are the outgoing flat vector
$L=(1,n)$, cost $O(r^{-2+0.04})$.  This uses
\eqref{far:point-first}, \eqref{far:projected-bounds}, and smooth inverse
composition.  The bad component of $X_g-L$ is $O(d_g)$, and the screen
is tangent to the sphere, so true outgoing and screen derivatives are
good derivatives for these estimates.  The exact lab gauge gives the
lowered wave contraction $O(r^{-2})$.  Thus the divergence comparison
for $r^{-2}d\operatorname{vol}_g$ and the deformation bounds, including
\Cref{fe:F16}, have their stated inputs.

\subsection{The far estimate}

The point bounds and the corrections above have two distinct uses.
The outgoing multiplier controls good derivatives and the corrected
cap flux; the laboratory-time multiplier controls bad derivatives.
Appendix~\ref{fe:section} proves both estimates, with explicit source
costs \eqref{far:source-costs} and \eqref{far:mixed-charge}.
Appendix~\ref{fs:section} estimates those costs by the ordered
commutations of the equation just verified.

\begin{proposition}[Far estimate]\label{far:estimate}
After giving the finitely many commuted energies fixed hierarchy weights,
there is $c_1>0$ such that
\begin{equation}\label{far:far-bound}
 \mathsf F\le o_{R_0}(1)(\mathsf F+\mathsf S_a)
       +C(R_0)\mathsf K_R
       +C_{\mathrm{fix}}\bigl(\text{input errors}
                                      +\mathsf N^{1+c_1}\bigr).
\end{equation}
The coefficient denoted by $o_{R_0}(1)$ tends to zero as $R_0\to\infty$,
independently of the upper compact matching radius $R$.  The input errors
vanish with $e_{\mathrm{in}}$, uniformly on bounded small bootstrap
ranges.
\end{proposition}

\begin{proof}
Apply the complete source estimate, Proposition~\ref{fs:source-estimate},
to the lab equation of Lemma~\ref{far:structure} and its $D$ derivative.
The lemma supplies its background symbol bounds, the additional $D$
gain on the defect and linear coefficients, and the differentiated
weak-null polarization.  The point bounds in
\eqref{far:point-h}--\eqref{far:news} supply the product estimates.
Lemma~\ref{far:boundary-input} supplies the full mixed-order,
physically nontrapping overlap and the initial symbol input.
These are the actual inputs of the six source estimates, rather than
bounds on an already global solution.

Two features explain the choice of levels.  A common principal
commutator places the differentiated coefficient on a lower field.
An extra bad laboratory-time derivative on $h$ can be replaced by
$\psi$ and good derivatives using $Dh=\psi$; the raw $h$ estimates
therefore use the prior strong level.  Conversely the differentiated equation is
\begin{equation}\label{far:strong-equation}
 g^{ij}\partial_{ij}\psi
   =Df_h-(Dg^{ij})\partial_{ij}h
      +2g^{ij}(\partial_is_s^d)\partial_{jd}h
      +g^{ij}(\partial_{ij}s_s^d)\partial_dh.
\end{equation}
Its seven commutations use $B^h$ indices through eight in the notation of
\eqref{fs:lists}: this means at most nine physical derivatives of $h$,
available at level $H$.  Its twice-null coefficient has the extra inverse radius because
\begin{equation}\label{far:strong-projection-identity}
 \begin{split}
 T_l^m\psi={}&(1-s_s\cdot n)T_l^{m+1}h
       +r^{-1}s_s^db_dZT_l^mh\\
 &+\sum_{j<m}c_{jm}(T_l^{m-j}s_s^d)\partial_dT_l^jh.
 \end{split}
\end{equation}
The differentiated gauge controls the first term without a stationary
background offset, since it is a derivative of $h$ itself.  The other
terms carry an inverse radius or an acceleration symbol.  The full
coefficient placements and transferred angular derivatives are checked
in Claims~1 and~2 of Proposition~\ref{fs:source-estimate}.

The pure defect uses a different estimate from field products.
After the $r^2$ source scaling it is $O(\widetilde a)$ for $h$
and $O(r^{-1}\widetilde a)$ for $\psi$.
The averaging estimates \eqref{fs:average-bounds} give squared shell
powers $p_i-1,p_i-A$ at the two $h$ levels and two additional inverse
powers at level $P$.  All are negative, as recorded in
\eqref{fs:defect-shell-margins}.  In particular its starting-annulus
cost has a small coefficient on $\mathsf S_a$, not a growing compact
coefficient.  The remaining source terms are handled by the strict
radial and time-weight margins in Claims~3--5; the leading strong
weak-null term uses the third current correction above.

Thus every correction, remainder, raw and mixed source budget has the
bound \eqref{fs:F15}.  The size-one terms respect the finite order:
level $P$ before the $h$ levels, lower total count, more lab times at
fixed count, and then fewer dilations.  The flat recurrence and its
nonnegative diagonal term are proved in Claim~3 of
Proposition~\ref{fe:energy-reduction}.  Its Claim~5 chooses fixed
hierarchy weights before increasing $R_0$; the source constants on
prior far energies are independent of both matching radii.
The two corrected multiplier estimates therefore give
\eqref{far:far-bound}.
\end{proof}

The background bound $B-\eta=O(r^{-1})$ may have a nonsmall fixed
constant.  The estimates above use its radial decay and differentiated
structure, so they do not prescribe the leading tail coefficient or
select a preferred Kerr member.

\subsection{Closing the exterior estimates on one base neighborhood}

The three inequalities now have the form
\begin{align}
 \mathsf K_R
 &\le CR^{-\kappa}(\mathsf F+\mathsf K_R)
    +C_{\mathrm{fix}}\bigl(e_{\mathrm{in}}+\mathsf N^{3/2}
                        +(e_{\mathrm{in}}\mathsf N)^{1/2}\bigr),
                                                   \label{far:compact-coupling}\\
 \mathsf F
 &\le\varepsilon_0(R_0)(\mathsf F+\mathsf S_a)
          +C(R_0)\mathsf K_R
          +C_{\mathrm{fix}}\bigl(\text{input errors}
                                      +\mathsf N^{1+c_1}\bigr),
                                                   \label{far:far-coupling}\\
 \mathsf S_a
 &\le C(R_0)\mathsf K_R+C_*\mathsf F
                   +C_{\mathrm{fix}}(e_{\mathrm{in}}+\mathsf N^2),
                                                   \label{far:sliced-coupling}
\end{align}
where $\kappa>0$, $\varepsilon_0(R_0)\to0$, and, crucially, $C_*$ is
independent of $R_0$.  The first inequality is \eqref{co2:high-bound}, the compact estimate with
mixed-order recovery.  Its acceleration and static quantities occur
only as small coefficient differences or products, not as linear
forcing.  The last inequality is \eqref{co2:sliced-acceleration}; its
radius-uniform far coefficient is proved in
Lemma~\ref{fs:uniform-coupling} by a single spatial adjoint pairing
and the strict margin $p_P-A=0.26$.  The compact-tail defect
comparison uses a cut $Y$ fixed before $R_0$.

\begin{proposition}[Simultaneous closure]\label{far:closure}
For each fixed strictly subextremal center and fixed computing cylinders,
there is a base smallness threshold depending only on the finite
$p_{10}$ input such that the compact, far, static, and modulation bounds
hold on every finite slab.  They yield global smooth vacuum evolution
on both computing cylinders.  Every fixed higher-order bound of
Proposition~\ref{co2:higher} then holds with datum- and order-dependent
constants; those higher seminorms need only be finite.
\end{proposition}

\begin{proof}
The common finite-slab norm includes the three compact levels with
mixed jets, the far levels including \eqref{far:corrected-norm},
$\mathsf S_a$, and the static norms of \eqref{co2:sliced-static}
on fixed reaches containing all compact error supports and deep
matching regions.  Include its extra nontrapped time orders and the
low norm $\mathsf l$ when needed.  Filtered feedback from the strong
$u$ level supplies the base parameter jets.  The static estimate
supplies the integrable coefficient difference near trapping; in the
compact high estimate it multiplies another bootstrap factor and
therefore creates no additional linear feedback term.

Choose $R_0$ so large that
$\varepsilon_0(R_0)(1+C_*)$ is as small as needed.  Substituting
\eqref{far:sliced-coupling} into \eqref{far:far-coupling} and absorbing
the resulting $\mathsf F$ term bounds $\mathsf F+\mathsf S_a$ by a fixed
$C(R_0)\mathsf K_R$ plus input and superlinear errors.  Now enlarge $R$
so that substitution in \eqref{far:compact-coupling} absorbs its compact
feedback.  Finally enlarge the plateau time $T_0$ after $R$ to obtain
the plateau-small terms used in that compact estimate.  This is the
order of \Cref{fs:coupled-inequalities}, and it has no radius cycle.

We spell out which earlier choices this order permits.  The inner
profile depth and finite Cauchy-join margins are fixed first.  The
slightly enlarged spacelike and outflow cylinders are therefore fixed
before their stationary gauges.  Measurement and gauge supports, the
high spatial cutoff $\Lambda$, the feedback filter, the small exponent
in the compact low transfer, and then the outer-current slack can all
be fixed independently of $R_0,R$.  Their polynomial frequency degrees
require only fixed differentiation, propagation, and testing orders.
Compact stationary energy gluing to the differential end, and the cut
$Y$ in the first sliced tail, are fixed before $R_0$ as well.
After choosing $R_0$, prior fixed compact reaches may include its
starting overlaps, with constants independent of $R$ and unchanged
exponents.  The original choice of $R_0$ also places the auxiliary
anchor exterior \eqref{far:anchor-domain} in the prescribed lab end.
Every constant needed for a radius-uniform absorption was fixed before
$R_0$; all other linear far costs were explicitly put on
$C(R_0)\mathsf K_R$.

The low norm and the second sliced norms cost only fixed linear
constants times the three absorbed norms, plus data and superlinear
errors.  Thus for some $c_2>0$,
\begin{equation}\label{far:total-closure}
 \mathsf N\le C_{\mathrm{fix}}
       \bigl(\mathcal R(e_{\mathrm{in}},\mathsf N)
                          +\mathsf N^{1+c_2}\bigr),
\end{equation}
where $\mathcal R(e_{\mathrm{in}},\mathsf N)\to0$ as
$e_{\mathrm{in}}\to0$, uniformly on bounded small bootstrap ranges.
This notation includes the mixed positive input powers in
\eqref{far:compact-coupling}.
Choose a bootstrap cap so small that the superlinear term is less than
one quarter of that cap, then choose the base input small enough that
the other term is less than one quarter.  The resulting strict
improvement is uniform in the terminal time.

For smooth finite-time solutions these norms vary continuously with the
slab endpoint.  For the cap suprema use the spacelike transversality of
\eqref{far:caps}, as in the construction of \Cref{fe:currents}.
Finite-time symbol bounds and strict radial convergence margins control
the tails.  Startup and
earlier bounded intervals have ordinary finite-time wave stability.
Since $A>1$, the acceleration bound keeps the parameters small, and the
spacelike, outflow, and wave-coefficient continuation conditions persist.
Proposition~\ref{st:local-evolution} now supplies causal local continuation
and constraint propagation, hence global smooth vacuum evolution on both
cylinders.  Its finite-time tame propagation of higher smooth orders
requires finiteness, not smallness, of those orders; no higher-order
trapping energy has been invoked to reach this conclusion.

All small initial norms used here are supplied by $p_{10}$ for the
fixed cylinder-slice construction.  The order-nine, eight, and seven
far commutations in \eqref{far:levels} carry one additional energy
derivative, exactly matching the compact field orders ten, nine, and
eight for $h,h,u$, respectively.  The higher-order induction is applied only after this base
closure.  Its later datum- and order-dependent starting times and
constants impose no further smallness on the common neighborhood.
\end{proof}
